\chapter{کاوش جریان‌های داده \texorpdfstring{\enparen{Data Streams}}{(Data Streams)}}
\label{ch:data-streams}

\section{مقدمه و مدل محاسباتی جریان داده}
در سیستم‌های سنتی پایگاه داده، فرض بر این است که داده‌ها ابتدا روی دیسک به صورت پایدار ذخیره می‌شوند و سپس پرس‌وجوهای تحلیلی روی داده‌های ایستا اجرا می‌گردند. اما در بسیاری از کاربردهای نوین کلان‌داده، داده‌ها به صورت پیوسته، با سرعتی سرسام‌آور و در حجم نامحدود تولید می‌شوند؛ به این نوع داده‌ها \textbf{جریان داده}\footnote{\lr{Data Stream}} گفته می‌شود.

نمونه‌های بارز جریان داده عبارت‌اند از:
\begin{itemize}
    \item کلیک‌های کاربران در وب‌سایت‌های پرترافیک مانند گوگل یا آمازون.
    \item داده‌های ارسالی از میلیون‌ها حسگر در شبکه‌های اینترنت اشیاء \enparen{IoT}.
    \item بسته‌های اطلاعاتی عبوری از روترهای شبکه اینترنت (نظیر پایش حملات DDoS و ناهنجاری‌های شبکه).
    \item پیام‌های منتشرشده در شبکه‌های اجتماعی نظیر توییتر.
\end{itemize}

\begin{definition}[مدل پردازش جریان داده]
در مدل محاسباتی جریان داده (شکل \ref{fig:stream-model}):
\begin{enumerate}
    \item داده‌ها به شکل دنباله‌ای بی‌پایان از تاپل‌ها وارد سیستم می‌شوند.
    \item نرخ ورود داده‌ها ممکن است بسیار فراتر از توان ذخیره‌سازی فوری روی دیسک باشد.
    \item سیستم پردازش جریان دارای یک \textbf{حافظه کاری}\footnote{\lr{Working Memory}} با ظرفیت محدود است.
    \item هر عنصر تنها \textbf{یک بار} (یا با تعداد دفعات بسیار محدود) قابل مشاهده است؛ اگر عنصری پردازش نشود، برای همیشه از دست می‌رود.
\end{enumerate}
\end{definition}

\begin{figure}[H]
\centering
\begin{tikzpicture}[
    scale=0.85, every node/.style={transform shape},
    box/.style={draw, rectangle, rounded corners, fill=blue!10, align=center, minimum height=1.1cm, minimum width=2.8cm, font=\small},
    arrow/.style={-{Latex[length=2.5mm]}, thick}
]
    \node[align=center] (in) at (-3, 0) {\rl{جریان داده ورودی}\\$\dots, s_3, s_2, s_1$};
    \node[box] (proc) at (1.5, 0) {\rl{پردازنده جریان}\\\rl{(الگوریتم‌های برخط)}};
    \node[box, fill=green!15] (mem) at (1.5, 2.2) {\rl{حافظه کاری محدود}\\\rl{(دسترسی سریع)}};
    \node[box, fill=orange!15] (store) at (1.5, -2.2) {\rl{حافظه ماندگار (دیسک)}\\\rl{(دسترسی کند، آرشیو)}};
    \node (out) at (6, 0) {\rl{پاسخ پرس‌وجوها}};
    
    \draw[arrow] (in) -- (proc);
    \draw[arrow] (proc) -- (out);
    \draw[<->, thick] (proc) -- (mem);
    \draw[arrow] (proc) -- (store);
\end{tikzpicture}
\caption{معماری سیستم مدیریت جریان داده: پردازش فوری با تکیه بر حافظه کاری محدود}
\label{fig:stream-model}
\end{figure}

پرس‌وجوها در جریان داده به دو دسته تقسیم می‌شوند:
\begin{itemize}
    \item \textbf{پرس‌وجوهای دائمی}\footnote{\lr{Standing Queries}}: پرس‌وجوهایی که همواره فعال هستند و به ازای هر عنصر ورودی، پاسخی به روز تولید می‌کنند (مثلاً: «اگر دمای هر حسگر از ۸۰ درجه بیشتر شد، هشدار بده»).
    \item \textbf{پرس‌وجوهای موردی}\footnote{\lr{Ad-hoc Queries}}: پرس‌وجوهایی که در یک لحظه خاص توسط تحلیل‌گر مطرح می‌شوند (مثلاً: «میانگین طول پیام‌های ارسال‌شده در ۱۰ دقیقه گذشته چقدر بوده است؟»).
\end{itemize}

\section{نمونه‌گیری از جریان‌های داده}
هنگامی که نمی‌توان تمام جریان را ذخیره کرد، نمونه‌گیری یک راهکار طبیعی است. دو مسئله اساسی در نمونه‌گیری وجود دارد:

\subsection{نمونه‌گیری با نسبت ثابت \texorpdfstring{\enparen{Fixed Proportion Sampling}}{(Fixed Proportion Sampling)}}
فرض کنید می‌خواهیم $10\%$ از پرس‌وجوهای یک موتور جستجو را نمونه‌برداری کنیم تا رفتار کاربران را بررسی نماییم.
اگر به ازای هر پرس‌وجو یک عدد تصادفی بین $0$ تا $9$ تولید کنیم و در صورت صفر بودن آن را برگزینیم، هر کاربر دارای تعداد پرس‌وجوهای متفاوتی در نمونه خواهد بود، اما اگر بخواهیم بفهمیم چه کسری از کاربران تنها یک بار پرس‌وجو کرده‌اند، این نمونه‌گیری نتایج را به شدت مخدوش می‌کند.

\textbf{راهکار اصولی}: نمونه‌گیری بر مبنای \textbf{کلید کاربر}. با استفاده از یک تابع درهم‌سازی:
\begin{equation}
h(\text{UserID}) \pmod{10} = 0
\end{equation}
اگر باقیمانده صفر شد، تمام پرس‌وجوهای آن کاربر ذخیره می‌شود و در غیر این صورت هیچ‌یک از پرس‌وجوهای او ذخیره نمی‌گردد. این کار نسبت نمونه $10\%$ را تضمین کرده و در عین حال ساختار کامل رفتار هر کاربر نمونه را حفظ می‌کند.

\subsection{نمونه‌گیری با اندازه ثابت: الگوریتم مخزنی \texorpdfstring{\enparen{Reservoir Sampling}}{(Reservoir Sampling)}}
فرض کنید می‌خواهیم در هر لحظه نمونه‌ای تصادفی و یکنواخت به اندازه دقیقاً $s$ از میان $n$ عنصر نخست جریان که تاکنون مشاهده شده‌اند در اختیار داشته باشیم، در حالی که مقدار نهایی $n$ از پیش معلوم نیست و حافظه ما فقط گنجایش $s$ عنصر را دارد.

\begin{algorithmbox}[الگوریتم نمونه‌گیری مخزنی \enparen{Reservoir Sampling}]
\begin{itemize}
    \item \textbf{ورودی}: جریان داده $x_1, x_2, \dots$ و ظرفیت مخزن $s$.
    \item \textbf{گام اول}: $s$ عنصر نخست جریان $(x_1, \dots, x_s)$ را مستقیماً در مخزن قرار دهید.
    \item \textbf{گام دوم}: برای هر عنصر ورودی بعدی $x_n$ (که $n > s$):
    \begin{enumerate}
        \item با احتمال $\frac{s}{n}$ عنصر $x_n$ را برای ورود به مخزن بپذیرید.
        \item در صورت پذیرش، یکی از $s$ عنصر موجود در مخزن را به صورت کاملاً تصادفی و یکنواخت (با احتمال $1/s$) انتخاب کرده و عنصر جدید را جایگزین آن نمایید.
        \item با احتمال $1 - \frac{s}{n}$ عنصر جدید را بدون تغییر مخزن دور بریزید.
    \end{enumerate}
\end{itemize}
\end{algorithmbox}

\begin{theorem}[صحت الگوریتم مخزنی]
پس از رسیدن $n$ عنصر از جریان، احتمال حضور هر عنصر $x_i$ $(1 \le i \le n)$ در مخزن دقیقاً برابر با $\frac{s}{n}$ است.
\end{theorem}

\textbf{اثبات با استقرای ریاضی:}
\begin{itemize}
    \item \textbf{پایه استقرا}: برای $n = s$، تمامی $s$ عنصر اول در مخزن قرار دارند و احتمال حضور هر یک $s/s = 1$ است که درست است.
    \item \textbf{فرض استقرا}: فرض کنیم حکم برای $n-1$ عنصر برقرار باشد؛ یعنی بعد از $n-1$ عنصر، هر عنصری با احتمال $\frac{s}{n-1}$ در مخزن قرار دارد.
    \item \textbf{گام استقرا}: حال عنصر $n$-ام می‌رسد:
    \begin{enumerate}
        \item برای خود عنصر $n$-ام، الگوریتم مستقیماً آن را با احتمال $\frac{s}{n}$ در مخزن قرار می‌دهد.
        \item برای هر عنصر قبلی $x_i$ $(i < n)$: این عنصر در مخزن باقی می‌ماند اگر و تنها اگر قبل از ورود عنصر $n$-ام در مخزن بوده باشد، و با ورود عنصر جدید حذف نشود:
        \begin{align}
        \Pr[x_i \text{ در مخزن بعد از گام } n] &= \Pr[x_i \text{ در مخزن در گام } n-1] \times \Pr[x_i \text{ حذف نشود}] \nonumber \\
        &= \left(\frac{s}{n-1}\right) \times \left( (1 - \frac{s}{n}) + \frac{s}{n} \left(1 - \frac{1}{s}\right) \right) \nonumber \\
        &= \left(\frac{s}{n-1}\right) \times \left( 1 - \frac{s}{n} + \frac{s}{n} - \frac{1}{n} \right) \nonumber \\
        &= \left(\frac{s}{n-1}\right) \times \left( \frac{n-1}{n} \right) = \frac{s}{n}
        \end{align}
    \end{enumerate}
    بنابراین حکم برای تمامی $n$ عنصر اثبات می‌شود.
\end{itemize}

\section{فیلتر کردن جریان‌های داده: فیلتر بلوم \texorpdfstring{\enparen{Bloom Filter}}{(Bloom Filter)}}
فرض کنید مجموعه‌ای از $m$ کلید مجاز یا نامطلوب $S$ داریم (مثلاً ۱ میلیارد آدرس ایمیل اسپم یا URLهای فیشینگ). جریان عظیمی از تاپل‌ها وارد می‌شود و می‌خواهیم در چند میکروثانیه بررسی کنیم آیا کلید ورودی در $S$ هست یا خیر. اندازه $S$ ممکن است به گیگابایت‌ها برسد و در حافظه کش سرور جا نگیرد.

\begin{definition}[فیلتر بلوم~\cite{bloom1970space}]
یک \textbf{فیلتر بلوم}\footnote{\lr{Bloom Filter}} یک ساختار داده احتمالی و بسیار فشرده برای آزمون عضویت در مجموعه است که:
\begin{itemize}
    \item دارای یک آرایه بیتی $B$ به طول $n$ (در ابتدا همگی $0$) است.
    \item از $k$ تابع درهم‌سازی مستقل $h_1, h_2, \dots, h_k$ استفاده می‌کند که بازه خروجی آنها $[0, n-1]$ است.
    \item \textbf{درج عضو $x \in S$}: برای تمام $i \in \{1, \dots, k\}$، بیت $B[h_i(x)]$ برابر با $1$ قرار داده می‌شود.
    \item \textbf{پرس‌وجوی عضویت $y$}: اگر برای تمامی $i \in \{1, \dots, k\}$، مقدار $B[h_i(y)] == 1$ باشد، پاسخ «بله» است؛ وگرنه «خیر».
\end{itemize}
\end{definition}

ویژگی حیاتی فیلتر بلوم:
\begin{itemize}
    \item \textbf{عدم وجود منفی کاذب \enparen{Zero False Negatives}}: اگر عنصری در $S$ باشد، بیت‌های مربوط به آن قطعاً ۱ شده‌اند و امکان ندارد پاسخ منفی دریافت کند.
    \item \textbf{احتمال مثبت کاذب \enparen{False Positives}}: ممکن است عنصری در $S$ نباشد، اما تصادفاً تمامی بیت‌های متناظر با توابع هش آن توسط سایر عناصر ۱ شده باشند.
\end{itemize}

\subsection{تحلیل احتمال خطای مثبت کاذب و مقدار بهینه \texorpdfstring{$k$}{k}}
فرض کنید $m$ عنصر در فیلتر با اندازه $n$ بیت و با $k$ تابع هش درج شده باشند:
\begin{enumerate}
    \item احتمال اینکه یک بیت خاص پس از یک بار اعمال هش همچنان $0$ باقی بماند برابر است با $1 - \frac{1}{n}$.
    \item پس از درج $m$ عنصر با $k$ تابع هش (مجموعاً $k \cdot m$ درج بیت)، احتمال اینکه یک بیت هنوز $0$ باشد:
    \begin{equation}
    p_0 = \left( 1 - \frac{1}{n} \right)^{km} \approx e^{-km/n}
    \end{equation}
    \item بنابراین، کسر بیت‌هایی از آرایه که مقدار $1$ به خود گرفته‌اند:
    \begin{equation}
    p_1 = 1 - e^{-km/n}
    \end{equation}
    \item احتمال اینکه برای یک عنصر ناموجود $y$، تمامی $k$ بیت هش‌شده برابر با ۱ باشند (نرخ مثبت کاذب):
    \begin{equation}
    P(\text{False Positive}) = (p_1)^k = \left( 1 - e^{-km/n} \right)^k
    \end{equation}
\end{enumerate}

با مشتق‌گیری نسبت به $k$ و مساوی صفر قرار دادن آن، \textbf{تعداد بهینه توابع درهم‌سازی} برای کمینه‌سازی خطای مثبت کاذب به دست می‌آید:
\begin{equation}
k_{\text{opt}} = \frac{n}{m} \ln 2 \approx 0.693 \times \frac{n}{m}
\end{equation}
و با جایگذاری این مقدار، کمترین احتمال مثبت کاذب ممکن برابر است با:
\begin{equation}
P_{\text{opt}} = \left( \frac{1}{2} \right)^k = (0.6185)^{n/m}
\end{equation}

\begin{example}[محاسبه پارامترهای فیلتر بلوم]
فرض کنید $m = 10^9$ کلید داریم و می‌خواهیم نرخ مثبت کاذب کمتر از $1\%$ $(0.01)$ باشد.
\begin{itemize}
    \item از رابطه $P \approx (0.5)^k \le 0.01 \implies 2^k \ge 100 \implies k \approx 7$.
    \item از رابطه $k = \frac{n}{m} \ln 2 \implies \frac{n}{m} = \frac{7}{\ln 2} \approx 10.1$ بیت به ازای هر کلید.
    \item بنابراین با اختصاص حدود ۱۰ بیت (تنها ۱.۲۵ بایت) به ازای هر کلید، خطای مثبت کاذب به زیر ۱ درصد می‌رسد!
\end{itemize}
\end{example}

\section{شمارش عناصر متمایز: الگوریتم فلژوله-مارتین \texorpdfstring{\enparen{Flajolet-Martin}}{(Flajolet-Martin)}}
مسئله: در یک جریان از شناسه‌ها، چه تعداد عنصر یکتا و متمایز وجود دارد؟
اگر بخواهیم شناسه‌ها را در جدول درهم‌سازی ذخیره کنیم، با افزایش تنوع عناصر با کمبود شدید حافظه مواجه می‌شویم. الگوریتم فلژوله-مارتین \enparen{FM}~\cite{flajolet1985probabilistic} تنها با مصرف حافظه‌ای در حد چند کلمه، این مقدار را تخمین می‌زند.

\begin{algorithmbox}[الگوریتم فلژوله-مارتین برای شمارش عناصر متمایز]
\begin{itemize}
    \item \textbf{تابع هش}: تابعی مانند $h(a)$ که عناصر را به رشته‌ای از بیت‌ها به طول $L$ نگاشت می‌کند (به طوری که $2^L > N$).
    \item \textbf{تابع صفرهای انتهایی $r(a)$}: تعداد صفرهای متوالی در سمت راست نمایش دودویی $h(a)$.
    (مثلاً اگر $h(a) = \dots 101000$ باشد، $r(a) = 3$ است).
    \item \textbf{نگه‌داری بیشینه صفرهای انتهایی}:
    \begin{equation}
    R = \max_{a \in \text{Stream}} r(a)
    \end{equation}
    \item \textbf{تخمین تعداد عناصر متمایز}:
    \begin{equation}
    \hat{D} = 2^R / \phi \quad \text{که در آن } \phi \approx 0.77351
    \end{equation}
\end{itemize}
\end{algorithmbox}

\textbf{بینش شهودی و تحلیل آماری:}
احتمال اینکه هش یک عنصر با حداقل $i$ صفر متوالی تمام شود برابر با $2^{-i}$ است.
اگر $D$ عنصر متمایز در جریان وجود داشته باشد، احتمال اینکه هیچ‌کدام از آن‌ها با $i$ صفر خاتمه نیابند:
\begin{equation}
(1 - 2^{-i})^D \approx e^{-D \cdot 2^{-i}}
\end{equation}
\begin{itemize}
    \item اگر $2^i \ll D$ باشد، احتمال مشاهده نکردن صفر انتهایی به سمت صفر میل می‌کند (قطعاً چنین عنصری دیده می‌شود).
    \item اگر $2^i \gg D$ باشد، این احتمال نزدیک به ۱ است (تقریباً محال است چنین عنصری دیده شود).
    \item نقطه گذار دقیقاً در $2^R \approx D$ اتفاق می‌افتد.
\end{itemize}

\subsection{تکنیک کاهش واریانس: میانگینِ میانه‌ها \texorpdfstring{\enparen{Median of Means}}{(Median of Means)}}
چون $2^R$ همواره توانی از دو است، واریانس این تخمین‌گر بالاست. برای رفع این نقیصه:
\begin{enumerate}
    \item از چندین تابع هش مستقل استفاده می‌کنیم.
    \item توابع هش را به گروه‌های کوچکی افراز می‌کنیم.
    \item در هر گروه، میانگین حسابی\footnote{\lr{Average}} مقادیر $2^R$ را محاسبه می‌کنیم تا واریانس کاهش یابد.
    \item میان گروه‌ها، میانه\footnote{\lr{Median}} مقادیر میانگین را برمی‌گزینیم تا مقادیر پرت و نامتعارف بی‌اثر شوند.
\end{enumerate}

\section{شمارش بیت‌ها در پنجره لغزان: الگوریتم DGIM}
فرض کنید جریانی از صفرها و یک‌ها داریم و می‌خواهیم در هر لحظه تعداد بیت‌های ۱ موجود در \textbf{پنجره لغزان}\footnote{\lr{Sliding Window}} به طول $N$ گذشته را بدانیم.
ذخیره کل $N$ بیت نیازمند $N$ بیت حافظه است که برای پنجره‌های بزرگ (مثلاً $N = 10^9$) بسیار پرهزینه است. الگوریتم DGIM (ابداع‌شده توسط داتار، گیونیس، ایندیک و موتوانی~\cite{datar2002maintaining}) این مسئله را در فضای $O(\log^2 N)$ و با خطای حداکثر $50\%$ حل می‌کند.

\begin{definition}[ساختار سطل‌های DGIM]
الگوریتم پنجره را به صورت مجموعه‌ای از \textbf{سطل‌ها}\footnote{\lr{Buckets}} خلاصه می‌کند:
\begin{enumerate}
    \item هر سطل شامل دو مقدار است: (۱) زمان انتهای سطل (یک برچسب زمانی $O(\log N)$ بیتی)؛ (۲) اندازه سطل (تعداد بیت‌های ۱ موجود در آن).
    \item اندازه هر سطل همواره \textbf{توانی از ۲} است $(1, 2, 4, 8, \dots)$.
    \item \textbf{قاعده ظرفیت}: به ازای هر اندازه سطل، \textbf{حداکثر دو سطل} (یک یا دو سطل) در پنجره مجاز است و هرگز سه سطل با اندازه یکسان نباید وجود داشته باشد.
    \item سطل‌ها هرگز هم‌پوشانی ندارند و سطل‌های با اندازه بزرگ‌تر، برچسب‌های زمانی قدیمی‌تری دارند.
\end{enumerate}
\end{definition}

\begin{algorithmbox}[به‌روزرسانی و پاسخ به پرس‌وجو در الگوریتم DGIM]
\textbf{عملیات به‌روزرسانی هنگام ورود بیت جدید:}
\begin{enumerate}
    \item اگر بیت جدید $0$ باشد، کاری انجام نمی‌شود (تنها برچسب زمان جاری پیش می‌رود).
    \item اگر بیت جدید $1$ باشد:
    \begin{itemize}
        \item سطلی با اندازه ۱ و با زمان جاری ایجاد می‌شود.
        \item اگر اکنون سه سطل با اندازه ۱ داشته باشیم، دو سطل قدیمی‌تر در هم ادغام شده و یک سطل با اندازه ۲ تشکیل می‌دهند.
        \item این ادغام به صورت زنجیره‌ای به سمت اندازه‌های بزرگ‌تر ادامه می‌یابد تا قانون «حداکثر دو سطل از هر اندازه» برقرار بماند.
    \end{itemize}
    \item هر سطلی که زمان انتهای آن قدیمی‌تر از $N$ واحد زمانی گذشته باشد، از انتهای پنجره حذف می‌گردد.
\end{enumerate}

\textbf{پاسخ به پرس‌وجوی «تعداد ۱ها در $k$ موقعیت اخیر» $(k \le N)$:}
\begin{itemize}
    \item تمامی سطل‌هایی را که کاملاً در بازه $k$ قرار دارند پیدا کرده و اندازه‌های آن‌ها را به طور کامل جمع بزنید.
    \item برای آخرین سطل که تا حدی با ابتدای پنجره تداخل دارد، دقیقاً \textbf{نیمی از اندازه آن} را به حاصل جمع اضافه کنید.
\end{itemize}
\end{algorithmbox}

\begin{theorem}[کران خطای الگوریتم DGIM]
حداکثر درصد خطای الگوریتم DGIM هرگز از $50\%$ تجاوز نمی‌کند.
\end{theorem}
\textbf{اثبات:}
فرض کنید آخرین سطل با اندازه $2^j$ دارای تداخل باشد. تخمین ما از این سطل برابر با $2^{j-1}$ است؛ لذا حداکثر خطای مطلق ممکن برابر با $2^{j-1} - 1$ است. از طرفی، چون قبل از این سطل حداقل یک سطل از هر یک از اندازه‌های $1, 2, 4, \dots, 2^{j-1}$ وجود دارد، مجموع واقعی ۱ها حداقل برابر با $1 + \sum_{i=0}^{j-1} 2^i = 2^j$ خواهد بود. بنابراین نسبت خطا:
\begin{equation}
\text{Error} \le \frac{2^{j-1}}{2^j} = 50\%
\end{equation}
(با افزایش سقف تعداد سطل‌ها از ۲ به $r$، می‌توان خطا را به دلخواه تا حد $1/r$ کاهش داد).

\section{تخمین گشتاورهای جریان: الگوریتم آلون-ماتیاس-سگدی \texorpdfstring{\enparen{AMS}}{(AMS)}}
در تحلیل آماری جریان، \textbf{گشتاور $k$-ام}\footnote{\lr{$k$-th Moment}} به صورت زیر تعریف می‌شود:
\begin{equation}
F_k = \sum_{v \in V} m_v^k
\end{equation}
که در آن $m_v$ فراوانی رخداد مقدار $v$ در جریان است:
\begin{itemize}
    \item \textbf{گشتاور صفرم $(F_0)$}: تعداد عناصر متمایز است (توسط الگوریتم فلژوله-مارتین حل شد).
    \item \textbf{گشتاور یکم $(F_1)$}: طول کل جریان $(N = \sum m_v)$.
    \item \textbf{گشتاور دوم $(F_2)$}: \textbf{عدد شگفتی \enparen{Surprise Number}} نام دارد و مقیاسی از میزان کجی (چولگی) یا ناهمگونی توزیع است. اگر همه عناصر فراوانی یکسان داشته باشند، $F_2$ کمینه است؛ و اگر یک عنصر کل جریان را تسخیر کند، $F_2$ بیشینه $(N^2)$ خواهد بود.
\end{itemize}

\begin{algorithmbox}[الگوریتم AMS برای محاسبه گشتاور دوم $F_2$]
الگوریتم Alon-Matias-Szegedy یک برآوردگر نااریب شگفت‌انگیز برای $F_2$ در فضای محدود است:
\begin{enumerate}
    \item یک موقعیت تصادفی $p$ در بازه زمانی $[1, N]$ انتخاب کرده و عنصر دیده‌شده در آن لحظه $(v)$ را ذخیره کنید.
    \item متغیری به نام $c$ تعریف کرده و تعداد دفعات مشاهده مجدد $v$ را از موقعیت $p$ تا انتهای جریان بشمارید (در ابتدا $c = 1$).
    \item در پایان، مقدار برآوردگر را به صورت زیر محاسبه نمایید:
    \begin{equation}
    X = N (2c - 1)
    \end{equation}
\end{enumerate}
\end{algorithmbox}

\begin{theorem}[نااریب بودن برآوردگر AMS]
امید ریاضی متغیر تصادفی $X$ در الگوریتم AMS دقیقاً برابر با گشتاور دوم واقعی $F_2$ است:
\begin{equation}
\mathbb{E}[X] = F_2
\end{equation}
\end{theorem}
\textbf{اثبات:}
فرض کنید عنصر $v$ با فراوانی کلی $m_v$ در جریان وجود دارد. احتمال اینکه موقعیت تصادفی $p$ بر روی یکی از این $m_v$ وقوع قرار گیرد برابر با $1/N$ است. اگر بر روی $i$-امین وقوع از آخر قرار گیرد، متغیر شمارش برابر با $i$ خواهد شد $(i \in \{1, \dots, m_v\})$. بنابراین:
\begin{equation}
\mathbb{E}[X] = \sum_{v \in V} \frac{1}{N} \sum_{i=1}^{m_v} N(2i - 1) = \sum_{v \in V} \sum_{i=1}^{m_v} (2i - 1)
\end{equation}
چون مجموع $m$ عدد فرد اول برابر با مربع کامل $m^2$ است $(\sum_{i=1}^{m} (2i - 1) = 1 + 3 + 5 + \dots + (2m-1) = m^2)$:
\begin{equation}
\mathbb{E}[X] = \sum_{v \in V} m_v^2 = F_2
\end{equation}
با نگه‌داری چندین نمونه متغیر تصادفی $X$ و محاسبه میانگین میانه‌های آن‌ها، واریانس تخمین به سرعت کاهش می‌یابد.

\section{پنجره‌های با کاهش نمایی \texorpdfstring{\enparen{Decaying Windows}}{(Decaying Windows)}}
در بسیاری از کاربردهای تحلیل جریان (نظیر روندهای توییتر یا معاملات بورس)، رخدادهای اخیر از اهمیت بسیار بالاتری نسبت به رویدادهای دوردست برخوردارند. در \textbf{پنجره با کاهش نمایی}\footnote{\lr{Decaying Window}}، وزن هر رویداد با گذشت زمان به صورت نمایی مستهلک می‌شود:
\begin{equation}
w(t) = e^{-\lambda (t_{\text{now}} - t)} \approx (1 - c)^{t_{\text{now}} - t}
\end{equation}
مزیت بی‌نظیر این رویکرد در محاسبات کلان‌داده این است که برای به‌روزرسانی مجموع اوزان در هنگام ورود یک عنصر جدید، نیازی به نگه‌داری هیچ تاریخچه‌ای نیست؛ تنها کافی است مقدار مجموع پیشین را در ضریب تخفیف $(1 - c)$ ضرب کرده و عدد ۱ را به آن بیفزاییم:
\begin{equation}
S_{\text{new}} = S_{\text{old}} \times (1 - c) + 1
\end{equation}
این فرمول در زمان و فضای ثابت $O(1)$ اجرا شده و بار محاسباتی سرورهای تحلیل بلادرنگ را به صفر نزدیک می‌کند.

\begin{summarybox}[جمع‌بندی فصل چهارم]
\begin{itemize}
    \item در جریان‌های داده، امکان ذخیره‌سازی کل تاریخچه وجود ندارد؛ پردازش باید یک‌گذره و با حافظه زیرخطی انجام شود.
    \item نمونه‌گیری مخزنی به انتخاب تصادفی یکنواخت از میان عناصر با اندازه نامعلوم جریان امکان می‌دهد.
    \item فیلتر بلوم با حداقل مصرف حافظه، عضویت در مجموعه‌ها را با صفر درصد منفی کاذب بررسی می‌کند.
    \item الگوریتم Flajolet-Martin شمارش عناصر متمایز را از روی بیشینه صفرهای متوالی درهم‌سازی تخمین می‌زند.
    \item الگوریتم DGIM شمارش عناصر پنجره لغزان را در فضای $O(\log^2 N)$ با خطای تضمین‌شده کنترل می‌نماید.
\end{itemize}
\end{summarybox}
